% GENERATED FILE. Edit canonical lessons, then run npm run generate:books.
% canonical-source-hash: fnv1a64:be38944b61041b41
% canonical-lessons: ZH-C437-shanqu, ZH-C437-pingyuan, ZH-C437-shamo, ZH-C437-gaoyuan, ZH-C437-yimin

\chapter{Mountains and Plains}
\label{ch:zh-mountains-and-plains}
\hlchaptermodality{437}

\begin{chapteropening}
\textbf{By the end of this chapter:} \emph{I can talk about a mountain area, a plain, a desert, a plateau and moving abroad.}

Complete the last lesson of chapter 437: I can talk about a mountain area, a plain, a desert, a plateau and moving abroad.
\end{chapteropening}

\section[\texorpdfstring{\zh{山区} (shānqū)}{山区 (shānqū)}]{\zh{山区} (shānqū) — mountain area}

\label{lesson:ZH-C437-shanqu}



\begin{quote}
Before the new one: say the Mandarin for a regional dialect, then the Mandarin for a regional accent.
\end{quote}

\subsection*{You'll want to know: \zh{山区}}
\textbf{\zh{山区}} — \emph{shānqū} — \textquotedblleft{}mountain area\textquotedblright{}.

\begin{quote}
\emph{wǒ nǎinai zhù zài shānqū} — My grandmother lives up in the mountains.
\end{quote}

\begin{cousinweb}
\textbf{\zh{山区}} (\emph{shānqū}) is built from \zh{山} 'a mountain' and \zh{区} 'an area', both inherited from Middle Chinese, \emph{srean} and \emph{khju} (Guangyun, in Baxter's transcription); the dictionary gives the second as the name of a marsh, and the sense 'an area' is not given there.
\end{cousinweb}

\begin{grammarlens}[title={What you've built}]
One more word for getting around.
\end{grammarlens}

\subsection*{Your turn}
\begin{itemize}
\raggedright
  \item \emph{Say it:} \emph{shānqū}
  \item \emph{Say it:} \emph{shānqū}, once more
  \item \emph{Say it:} \emph{wǒ nǎinai zhù zài shānqū} — My grandmother lives up in the mountains.
  \item \emph{Recall:} say the Mandarin for a regional dialect, then the Mandarin for a regional accent, then say \emph{shānqū} again
\end{itemize}

\begin{tcolorbox}[breakable,colback=teal!4,colframe=teal!35!black,title={Before you move on}]
What does \zh{山区} mean? (\textquotedblleft{}Mountain area\textquotedblright{}.) Say it once more.
\end{tcolorbox}
\section[píngyuán]{píngyuán — plain; flat open land}

\label{lesson:ZH-C437-pingyuan}



\begin{quote}
Before the new one: say the Mandarin for a regional accent, then the Mandarin for a mountain area.
\end{quote}

\subsection*{You'll want to know: píngyuán}
\emph{píngyuán} — \textquotedblleft{}plain; flat open land\textquotedblright{}.

\begin{quote}
\emph{běifāng de píngyuán hěn dà} — The plains in the north are huge.
\end{quote}

Its characters come later; for now it is a word for the ear.

\begin{cousinweb}
\emph{píngyuán} is built from \emph{píng} 'level' and \emph{yuán} 'broad flat land', both inherited from Middle Chinese, \emph{bjaeng} and \emph{ngjwon} (Guangyun, in Baxter's transcription), which glosses the first 'right' and 'at peace' and the second, after an older dictionary, as broad flat land; the sense 'level' for the first is not given there.
\end{cousinweb}

\begin{grammarlens}[title={What you've built}]
One more word for getting around.
\end{grammarlens}

\subsection*{Your turn}
\begin{itemize}
\raggedright
  \item \emph{Say it:} \emph{píngyuán}
  \item \emph{Say it:} \emph{píngyuán}, once more
  \item \emph{Say it:} \emph{běifāng de píngyuán hěn dà} — The plains in the north are huge.
  \item \emph{Recall:} say the Mandarin for a regional accent, then the Mandarin for a mountain area, then say \emph{píngyuán} again
\end{itemize}

\begin{tcolorbox}[breakable,colback=teal!4,colframe=teal!35!black,title={Before you move on}]
What does píngyuán mean? (\textquotedblleft{}Plain; flat open land\textquotedblright{}.) Say it once more.
\end{tcolorbox}
\section[shāmò]{shāmò — desert}

\label{lesson:ZH-C437-shamo}



\begin{quote}
Before the new one: say the Mandarin for a mountain area, then the Mandarin for flat open land.
\end{quote}

\subsection*{You'll want to know: shāmò}
\emph{shāmò} — \textquotedblleft{}desert\textquotedblright{}.

\begin{quote}
\emph{míngtiān wǒmen zuò chē qù shāmò} — Tomorrow we're driving out to the desert.
\end{quote}

Its characters come later; for now it is a word for the ear.

\begin{cousinweb}
\emph{shāmò} is built from \emph{shā} 'sand' and \emph{mò} 'a waste', both inherited from Middle Chinese, \emph{srae} and \emph{mak} (Guangyun, in Baxter's transcription), which glosses the first as stones scattered by water and gives the pair itself under the second.
\end{cousinweb}

\begin{grammarlens}[title={What you've built}]
One more word for getting around.
\end{grammarlens}

\subsection*{Your turn}
\begin{itemize}
\raggedright
  \item \emph{Say it:} \emph{shāmò}
  \item \emph{Say it:} \emph{shāmò}, once more
  \item \emph{Say it:} \emph{míngtiān wǒmen zuò chē qù shāmò} — Tomorrow we're driving out to the desert.
  \item \emph{Recall:} say the Mandarin for a mountain area, then the Mandarin for flat open land, then say \emph{shāmò} again
\end{itemize}

\begin{tcolorbox}[breakable,colback=teal!4,colframe=teal!35!black,title={Before you move on}]
What does shāmò mean? (\textquotedblleft{}Desert\textquotedblright{}.) Say it once more.
\end{tcolorbox}
\section[gāoyuán]{gāoyuán — plateau; highland}

\label{lesson:ZH-C437-gaoyuan}



\begin{quote}
Before the new one: say the Mandarin for flat open land, then the Mandarin for a desert.
\end{quote}

\subsection*{You'll want to know: gāoyuán}
\emph{gāoyuán} — \textquotedblleft{}plateau; highland\textquotedblright{}.

\begin{quote}
\emph{xīnán yǒu hěn duō gāoyuán} — There are a lot of high plateaus in the southwest.
\end{quote}

Its characters come later; for now it is a word for the ear.

\begin{cousinweb}
\emph{gāoyuán} is built from \zh{高} 'high' and \emph{yuán} 'broad flat land', both inherited from Middle Chinese, \emph{kaw} and \emph{ngjwon} (Guangyun, in Baxter's transcription), which glosses the first 'high' and 'lofty'.
\end{cousinweb}

\begin{grammarlens}[title={What you've built}]
One more word for getting around.
\end{grammarlens}

\subsection*{Your turn}
\begin{itemize}
\raggedright
  \item \emph{Say it:} \emph{gāoyuán}
  \item \emph{Say it:} \emph{gāoyuán}, once more
  \item \emph{Say it:} \emph{xīnán yǒu hěn duō gāoyuán} — There are a lot of high plateaus in the southwest.
  \item \emph{Recall:} say the Mandarin for flat open land, then the Mandarin for a desert, then say \emph{gāoyuán} again
\end{itemize}

\begin{tcolorbox}[breakable,colback=teal!4,colframe=teal!35!black,title={Before you move on}]
What does gāoyuán mean? (\textquotedblleft{}Plateau; highland\textquotedblright{}.) Say it once more.
\end{tcolorbox}
\section[yímín]{yímín — to emigrate; to settle abroad}

\label{lesson:ZH-C437-yimin}



\begin{quote}
Before the new one: say the Mandarin for a desert, then the Mandarin for a high plateau.
\end{quote}

\subsection*{You'll want to know: yímín}
\emph{yímín} — \textquotedblleft{}to emigrate; to settle abroad\textquotedblright{}.

\begin{quote}
\emph{tāmen quánjiā yímín qù le guówài} — Their whole family emigrated abroad.
\end{quote}

Its characters come later; for now it is a word for the ear.

\begin{cousinweb}
\emph{yímín} is built from \emph{yí} 'to move' and \emph{mín} 'the people', both inherited from Middle Chinese, \emph{ye} and \emph{mjin} (Guangyun, in Baxter's transcription), which glosses the first 'to shift' and the second 'the common people'.
\end{cousinweb}

\begin{grammarlens}[title={What you've built}]
One more word for getting around.
\end{grammarlens}

\subsection*{Your turn}
\begin{itemize}
\raggedright
  \item \emph{Say it:} \emph{yímín}
  \item \emph{Say it:} \emph{yímín}, once more
  \item \emph{Say it:} \emph{tāmen quánjiā yímín qù le guówài} — Their whole family emigrated abroad.
  \item \emph{Recall:} say the Mandarin for a desert, then the Mandarin for a high plateau, then say \emph{yímín} again
\end{itemize}

\begin{tcolorbox}[breakable,colback=teal!4,colframe=teal!35!black,title={Before you move on}]
What does yímín mean? (\textquotedblleft{}To emigrate; to settle abroad\textquotedblright{}.) Say it once more.
\end{tcolorbox}
